% !TeX root = ../thesis.tex
\chapter{Discussion}
\label{chp:discussion}

\section{Interpretation of results}
The central question behind this thesis is whether a single tool can bridge the gap between classroom exercise tools and controlled research software. The results from the classroom deployment with 107 participants across four conditions and four training days provide a first answer.

\subsection*{Did the tool work as a research instrument?}
On the research side, RekenRangers demonstrated that it can support a real experiment from start to finish. Two researchers with no programming background configured a 2$\times$2 factorial study, deployed it across seven schools, ran it over multiple days, monitored participant progress through the dashboard, and exported structured data ready for analysis. There were problems along the way (the duplicate submission bug, a server downtime incident, a demo that was too long) but none of these resulted in permanent data loss or forced the study to be abandoned. The final results of the study itself are not yet available, as Eline T'joen and Kiara Willems will complete their master's thesis in the next academic year. But the fact that the tool carried a study of this scale and complexity is itself a meaningful result. It shows that the gap identified in Chapters \ref{chp:background} and \ref{chp:exercise-tools} is not just a theoretical concern but one that can be addressed in practice.

\subsection*{Did the engagement system deliver?}
The engagement data tells a clearer story. Children spent 98.5\% of their earned coins on puzzle pieces, and 99\% completed at least one puzzle, all without any instruction or incentive to do so. The variation in purchasing strategies (some children going for cheap puzzles to rescue as many animals as possible, others saving for expensive ones) suggests that children were making deliberate choices rather than clicking through the shop mechanically. From the qualitative feedback, Eline reported that children were visibly enthusiastic about the puzzles and the safari theme. What the data cannot show is whether this engagement translated into better arithmetic performance or more sustained accuracy over the four training days. Demonstrating a causal link between gamification and learning outcomes would require a dedicated study with engagement as the independent variable, which was outside the scope of this thesis.

\subsection*{Did the adaptivity system perform as designed?}
The adaptive difficulty system produced mixed results. On the positive side, it clearly differentiated between children of different ability levels: final Elo scores ranged from 500 to 1000, and the trajectory data shows the system responding to each child individually. Children in the adaptive conditions received exercises closer to their own level (mean difficulty of 2.54 compared to the fixed 2.00 in non-adaptive conditions), and the response time data confirms that the scoring function pushed children toward exercises where they needed close to the full time available. These are signs that the system is doing what it was designed to do: keeping children in a zone where the exercises are challenging but not impossible. On the other hand, the system did not converge on its 75\% accuracy target. The mean accuracy of 68\% in the adaptive conditions indicates that the current calibration is too aggressive. The most likely causes are the fixed expected score of 0.75 and the scoring function \ref{for:scoring-function}, both of which were set based on limited pilot data. A more refined approach would be to assign an Elo rating to each individual exercise rather than working with fixed difficulty categories, following the item-level model used in Rekentuin \cite{klinkenberg2011computer}. This would allow the system to learn which specific exercises are hard and which are easy, rather than relying on a predefined classification. The ceiling effect, with 14 children reaching the maximum Elo of 1000, points in a similar direction: the current exercise pool does not extend far enough at the top end, and either harder exercises or a wider rating scale would be needed to keep the system informative for the strongest children. These are calibration and content problems rather than architectural ones. The adaptive framework itself is sound and can be improved incrementally.\\ \\
The response time difference between the adaptive and non-adaptive groups is worth a brief note. The adaptive group clustered toward longer response times, while the non-adaptive group was spread more uniformly. This is a direct consequence of the scoring function, whose equilibrium point sits at the time limit. The pattern confirms that the system behaves as designed, but it also raises the question of whether pushing children to work at the edge of their time limit is pedagogically desirable. A scoring function with a more generous equilibrium might produce a more comfortable experience without sacrificing the adaptive mechanism.



\section{Limitations}
\label{sec:limitations}

\section{Collaborating across disciplines: reflections for a computer scientist}

